\documentclass[10pt,a4paper]{amsart}
\usepackage{amsmath,amssymb,amsthm,graphicx}
% Retained minimal preamble for the Stacks Project chapter flat.tex (More on Flatness).
% Verbatim excerpt of the chapter preamble: Xy-pic setup (native preamble lines 14-18) and the native
% theorem declarations with the mathematical macros (native preamble lines 154-217). Copied mechanically;
% nothing added, removed or redefined. The wrapper class and the amsmath/amssymb/amsthm/graphicx base
% packages are supplied by the assembler, so the class/externaldocument part of the native preamble is omitted.
% For commutative diagrams we use Xy-pic
\usepackage[all]{xy}

% We use 2cell for 2-commutative diagrams.
\xyoption{2cell}
% Theorem environments.
%
\theoremstyle{plain}
\newtheorem{theorem}[subsection]{Theorem}
\newtheorem{proposition}[subsection]{Proposition}
\newtheorem{lemma}[subsection]{Lemma}

\theoremstyle{definition}
\newtheorem{definition}[subsection]{Definition}
\newtheorem{example}[subsection]{Example}
\newtheorem{exercise}[subsection]{Exercise}
\newtheorem{situation}[subsection]{Situation}

\theoremstyle{remark}
\newtheorem{remark}[subsection]{Remark}
\newtheorem{remarks}[subsection]{Remarks}

\numberwithin{equation}{subsection}

% Macros
%
\def\lim{\mathop{\mathrm{lim}}\nolimits}
\def\colim{\mathop{\mathrm{colim}}\nolimits}
\def\Spec{\mathop{\mathrm{Spec}}}
\def\Hom{\mathop{\mathrm{Hom}}\nolimits}
\def\Ext{\mathop{\mathrm{Ext}}\nolimits}
\def\SheafHom{\mathop{\mathcal{H}\!\mathit{om}}\nolimits}
\def\SheafExt{\mathop{\mathcal{E}\!\mathit{xt}}\nolimits}
\def\Sch{\mathit{Sch}}
\def\Mor{\mathop{\mathrm{Mor}}\nolimits}
\def\Ob{\mathop{\mathrm{Ob}}\nolimits}
\def\Sh{\mathop{\mathit{Sh}}\nolimits}
\def\NL{\mathop{N\!L}\nolimits}
\def\CH{\mathop{\mathrm{CH}}\nolimits}
\def\proetale{{pro\text{-}\acute{e}tale}}
\def\etale{{\acute{e}tale}}
\def\QCoh{\mathit{QCoh}}
\def\Ker{\mathop{\mathrm{Ker}}}
\def\Im{\mathop{\mathrm{Im}}}
\def\Coker{\mathop{\mathrm{Coker}}}
\def\Coim{\mathop{\mathrm{Coim}}}

% Boxtimes
%
\DeclareMathSymbol{\boxtimes}{\mathbin}{AMSa}{"02}

%
% Macros for moduli stacks/spaces
%
\def\QCohstack{\mathcal{QC}\!\mathit{oh}}
\def\Cohstack{\mathcal{C}\!\mathit{oh}}
\def\Spacesstack{\mathcal{S}\!\mathit{paces}}
\def\Quotfunctor{\mathrm{Quot}}
\def\Hilbfunctor{\mathrm{Hilb}}
\def\Curvesstack{\mathcal{C}\!\mathit{urves}}
\def\Polarizedstack{\mathcal{P}\!\mathit{olarized}}
\def\Complexesstack{\mathcal{C}\!\mathit{omplexes}}
% \Pic is the operator that assigns to X its picard group, usage \Pic(X)
% \Picardstack_{X/B} denotes the Picard stack of X over B
% \Picardfunctor_{X/B} denotes the Picard functor of X over B
\def\Pic{\mathop{\mathrm{Pic}}\nolimits}
\def\Picardstack{\mathcal{P}\!\mathit{ic}}
\def\Picardfunctor{\mathrm{Pic}}
\def\Deformationcategory{\mathcal{D}\!\mathit{ef}}

% --- appended 2026-10-12 (Chief): the project own definitions of the reference/slogan/history
% display environments, copied mechanically from the shared Stacks preamble lines 8-12
% (\usepackage{verbatim} plus the three \newenvironment lines, which the chapter preamble keeps
% verbatim from the project preamble). Nothing added, removed or redefined: these environments
% carry no mathematical content and the project itself typesets them as comments. Without these
% three lines a counted statement whose own text uses them fails with Environment slogan undefined.
\usepackage{verbatim}
\newenvironment{reference}{\comment}{\endcomment}
\newenvironment{slogan}{\comment}{\endcomment}
\newenvironment{history}{\comment}{\endcomment}
\usepackage{cleveref}
\crefname{originallemma}{original lemma}{original lemmas}
\Crefname{originallemma}{Original lemma}{Original lemmas}
\makeatletter
\newlabel{more-algebra-lemma-fitting-ideal-basics}{{07ZA}{source}{}{sourcepointer-1}{}}
\newlabel{more-algebra-lemma-fitting-ideal-basics@cref}{{[originallemma][0][]07ZA}{[1][0][]source}{}{}{}}
\makeatother
\crefname{originallemma}{original lemma}{original lemmas}
\Crefname{originallemma}{Original lemma}{Original lemmas}
\makeatletter
\newlabel{algebra-lemma-map-between-etale}{{00U7}{source}{}{sourcepointer-2}{}}
\newlabel{algebra-lemma-map-between-etale@cref}{{[originallemma][0][]00U7}{[1][0][]source}{}{}{}}
\makeatother
\crefname{originallemma}{original lemma}{original lemmas}
\Crefname{originallemma}{Original lemma}{Original lemmas}
\makeatletter
\newlabel{algebra-lemma-surjective-flat-finitely-presented}{{00U8}{source}{}{sourcepointer-3}{}}
\newlabel{algebra-lemma-surjective-flat-finitely-presented@cref}{{[originallemma][0][]00U8}{[1][0][]source}{}{}{}}
\makeatother
\title{The zeroth Fitting ideal of a quotient along a finite locally free ring map factors as an ideal of the base times an ideal after localising}
\author{The Stacks Project Authors}
\date{stacks-project:flat.tex:local-snapshot-20261011; GNU FDL 1.2 or later}
\begin{document}
\maketitle
\noindent\textit{Editorial note.} All delivered mathematical text is the author's own, copied from the pinned Stacks Project chapter file flat.tex; no statement or proof step was rewritten, reordered or removed.\par\medskip
\noindent\textit{Editorial note.} The preamble is the chapter's own retained minimal preamble (sources/stacks-flat/retained-preamble.tex), a verbatim excerpt of the native Xy-pic setup and of the native theorem declarations and mathematical macros, followed by the project own definitions of the reference, slogan and history display environments (verbatim package plus three newenvironment lines) that the retained file carries at its end and that the project itself typesets as comments. No macro was invented or redefined, and the wrapper class and the amsmath/amssymb/amsthm/graphicx base packages are supplied by the assembler.\par\medskip
\noindent\textit{Editorial note.} The counted statement and proof contain no cross reference, no citation and no included file, so the delivered body is self contained.\par\medskip
\noindent\textit{Editorial note.} Canonical location of the counted proposition: Stacks tag 0812 (sources/stacks-flat/pinned-original-source.tex line 8849).\par\medskip
\noindent\textit{Allowed standard prerequisites (editorial).}\begin{itemize}
\item Basic properties of Fitting ideals, in particular their behaviour under base change. More on Algebra, Lemma labelled more-algebra-lemma-fitting-ideal-basics (Stacks tag 07ZA), invoked at proof lines 8858-8860. Supplies that IC is the Fitting ideal of C/J tensor\_A C, which is the base change step that makes the splitting argument possible.
\item A map between etale algebras is etale, and a surjective flat finitely presented algebra map splits off a direct summand. Algebra, Lemmas labelled algebra-lemma-map-between-etale (Stacks tag 00U7) at proof lines 8864-8865 and algebra-lemma-surjective-flat-finitely-presented (Stacks tag 00U8) at proof lines 8866-8868. Supply that the multiplication map is etale, hence flat, and that the resulting surjection C tensor\_A C\_h -> C\_h splits, giving the direct sum decomposition used to factor the ideal.
\end{itemize}\medskip
\ifdefined\hypertarget\hypertarget{sourcepointer-1}{}\fi
\noindent\textit{Source reference.} Original lemma 07ZA: sources/stacks-chow/tags.txt, source lines 10140--10140. Cross-chapter standard prerequisite located in the project tag list staged in the local snapshot; invoked at proof line 8859 of the counted proof. It is retained only as a navigation pointer; every load-bearing step is inside the delivered span.\par\medskip
\ifdefined\hypertarget\hypertarget{sourcepointer-2}{}\fi
\noindent\textit{Source reference.} Original lemma 00U7: sources/stacks-chow/tags.txt, source lines 1101--1101. Cross-chapter standard prerequisite located in the project tag list staged in the local snapshot; invoked at proof line 8865 of the counted proof. It is retained only as a navigation pointer; every load-bearing step is inside the delivered span.\par\medskip
\ifdefined\hypertarget\hypertarget{sourcepointer-3}{}\fi
\noindent\textit{Source reference.} Original lemma 00U8: sources/stacks-chow/tags.txt, source lines 1102--1102. Cross-chapter standard prerequisite located in the project tag list staged in the local snapshot; invoked at proof line 8868 of the counted proof. It is retained only as a navigation pointer; every load-bearing step is inside the delivered span.\par\medskip
\begin{lemma}
\label{lemma-trick-fitting-ideal}
Let $A \to C$ be a finite locally free ring map of rank $d$.
Let $h \in C$ be an element such that $C_h$ is \'etale over $A$.
Let $J \subset C$ be an ideal. Set $I = \text{Fit}_0(C/J)$ where we
think of $C/J$ as a finite $A$-module. Then $IC_h = JJ'$ for some ideal
$J' \subset C_h$. If $J$ is finitely generated so are $I$ and $J'$.
\end{lemma}

\begin{proof}
We will use basic properties of Fitting ideals, see
More on Algebra, Lemma \ref{more-algebra-lemma-fitting-ideal-basics}.
Then $IC$ is the Fitting ideal of $C/J \otimes_A C$.
Note that $C \to C \otimes_A C$, $c \mapsto 1 \otimes c$ has a section
(the multiplication map). By assumption $C \to C \otimes_A C$ is \'etale
at every prime in the image of $\Spec(C_h)$ under this section.
Hence the multiplication map $C \otimes_A C_h \to C_h$ is \'etale
in particular flat, see Algebra, Lemma \ref{algebra-lemma-map-between-etale}.
Hence there exists a $C_h$-algebra such that
$C \otimes_A C_h \cong C_h \oplus C'$ as $C_h$-algebras, see
Algebra, Lemma \ref{algebra-lemma-surjective-flat-finitely-presented}.
Thus $(C/J) \otimes_A C_h \cong (C_h/J_h) \oplus C'/I'$ as
$C_h$-modules for some ideal $I' \subset C'$.
Hence $IC_h = JJ'$ with $J' = \text{Fit}_0(C'/I')$ where we view
$C'/J'$ as a $C_h$-module.
\end{proof}

\end{document}
